\section{视觉算法功能设计}
在概念层面，视觉（导航）系统可拆分为七个相互解耦的算法功能模块：激光惯性里程计、高频位姿插值、定位健康监测、坐标变换与位姿转发、飞行指令生成、统一安全管线、飞行状态监测与安全夺权。各模块以独立 ROS 2 节点或独立函数单元的形式运行，通过话题解耦，任一模块失效均不会导致整体系统崩溃，从架构上保证了系统的鲁棒性与可调试性。

\subsection{激光惯性里程计（FAST-LIO）}
激光惯性里程计负责将机载 3D 激光雷达点云与 IMU 数据进行紧耦合融合，输出全局一致的里程计位姿。选用 FAST-LIO 方案：其基于迭代扩展卡尔曼滤波（iEKF），以 IMU 高频预测、雷达点云低频校正，兼具精度高、计算量小、对退化环境鲁棒的特点，适合机载算力受限与快速运动场景。里程计同时发布定位状态位标志（\lstinline{/localization_status}，bit0=位姿可用、bit1=导航可用），供控制与监控节点判断定位是否可信。

\subsection{高频位姿插值}
由于雷达里程计帧率较低（约$\qty{10}{\Hz}$），直接用于$\qty{50}{\Hz}$的控制回路会引入明显延迟。为此利用 IMU 数据对里程计进行插值，输出高频（$\qty{100}{\Hz}$量级）的平滑位姿，显著降低控制延迟与位姿抖动，是保证飞行控制品质的关键环节。

\subsection{定位健康监测}
为保证“永不坠机”，在位姿转发链路中内置健康监测逻辑：对里程计位置做数值有限性校验与坐标越界校验，对连续帧做跳变检测与帧超时检测。一旦发现异常即判定定位不可信并\textbf{锁存}该状态，停止向飞控转发视觉位姿，使飞控因位姿缺失自动退出视觉依赖，从而规避漂移位姿导致的失控风险。

\subsection{坐标变换与位姿转发}
将雷达里程计坐标系（机体 FLU、世界 ENU）变换为飞控期望的坐标系（MAVROS 内部约定 ENU、PX4 使用 NED），同步对速度、协方差矩阵做旋转映射，并通过 \lstinline{mavros} 转发至飞控。同时以 \lstinline{MAV_STATE_ACTIVE} / \lstinline{MAV_STATE_FLIGHT_TERMINATION} 状态字向飞控报告视觉定位进程的健康状况。

\subsection{飞行指令生成}
导航控制节点根据协议指令（\lstinline{std_msgs/UInt8MultiArray}）或遥控器拨杆选择飞行模式并生成控制指令，支持三种模式：
\begin{itemize}
    \item \textbf{方向模式（Direction）}：按前后左右/转向位标志输出恒定方向速度，并叠加直线保持、保高与保向修正；
    \item \textbf{航点模式（DirectSetpoint）}：对目标航点进行比例导航逼近，近场自动减速，到达后自动悬停并锁定航向；
    \item \textbf{反反无人机模式（AntiAntiDrone）}：在锚点悬停后于上下两个高度目标间周期切换，规避对空打击。
\end{itemize}
指令来源支持外部话题（协议）与遥控器预设航点（RC）双源，并可在运行中通过拨杆动态切换。

\subsection{统一安全管线}
无论指令来源于何处，所有输出指令均依次经过统一安全管线处理后再下发：有限性校验 $\rightarrow$ 速度限幅 $\rightarrow$ 加速度（斜坡）约束 $\rightarrow$ 地理围栏与最低高度约束 $\rightarrow$ 发布。该管线将机体运动能力与场地边界固化为软件约束，从根本上消除速度阶跃、超速与越界风险。

\subsection{飞行状态监测与安全夺权}
飞行状态监测节点将飞控连接、位姿可用、导航可用三级判定汇总为飞行状态（\lstinline{/drone/flight_status}）上报电控；遥控器摇杆监测节点在 OFFBOARD 模式下检测到操作手拨动摇杆时，主动请求飞控切换至 ALTCTL 模式，实现控制权一键交还，作为自动飞行失效时的最终人工兜底。
